\begin{proof}[Proof of Theorem~\ref{thm: main}.]
There exists an $m$-neighborhood $V_0\subset M$
on which there exist smooth functions~$y^1,\ldots,y^k$
vanishing at~$m$
so that the leaves of $\d y^1 =\cdots =\d y^k = 0$
are intersections of the leaves of~$\cL$ with~$V_0$
as well as functions~$p_2,\ldots,p_k$, also vanishing
at~$m$, and a nonvanishing function~$a$ so that
\[
V_0^*\omega
= a\bigl(\d y^1 + p_2\,\d y^2 +\cdots + p_k\,\d y^k\bigr)
\]
By reversing the signs of~$a$ and the~$y^i$,
if necessary, one can assume that~$a(m)>0$.
Let~$W\subset T_mM$ be the tangent to the leaf of~$\cL$
that passes through~$m$, so that $W$ is the common kernel
of the~$\d y^i$ evaluated at~$m$.

Set~$\bar\omega = a^{-1}\omega$ and note that,
since~$\d\bar\omega \equiv a^{-1}\,\d\omega \bmod \omega$,
it follows that $\cL$ is also $\bar\omega$-Legendrian.
Moreover, since
\[
\Ss^\nabla\bar\omega = \d\big(a^{-1}\big)\circ\omega + a^{-1} \Ss^\nabla\omega,
\]
it follows that the tangent spaces of~$\cL$
(which, of course, satisfy $\omega=0$)
are also $\nabla$-positive for~$\bar\omega$.
Since~$\omega = a \bar\omega$ and $a>0$,
finding the desired convex representation for $\bar\omega$
will also yield one for~$\omega$.
Thus, it suffices to prove the theorem with~$\bar\omega$
in the place of~$\omega$, i.e., to assume that~$a=1$,
so I will do that from now on. Thus,
\[
V_0^*\omega = \d y^1 + p_2\,\d y^2 +\cdots + p_k\,\d y^k.
\]
Since~$\omega\w(\d\omega)^{k-1}\not=0$, the functions
$y^1,\ldots,y^k$, $p_2,\ldots,p_k$ have linearly independent differentials
at~$m$.

Restricting to~$V_0$, i.e., setting $M=V_0$, one has
\[
\Ss^\nabla\omega = H\big(y^1\big) + \d p_2\circ \d y^2 + \cdots + \d p_k\circ \d y^k
 + p_2 H\big(y^2\big) + \cdots + p_k H\big(y^k\big).
\]
Since the~$p_j$ vanish at~$m$,
it follows that, when restricted to~$W\subset T_mM$,
the two quadratic forms~$H\big(y^1\big)$ and $\Ss^\nabla\omega$ are equal.
Thus~$H\big(y^1\big)$ is positive definite on~$W$,
and so there is a~constant~$c>0$
so that~\smash{$H\big(y^1\big) + c \big(\d y^2\big)^2 + \cdots + c \big(\d y^k\big)^2$}
is positive definite on~$K_m = \{v\in T_mM\ \vrule\ \d y^1(v) = 0\}$.
Writing
\[
\omega = \d\bigl(y^1 + \tfrac12c \big(y^2\big)^2 + \cdots + \tfrac12c \big(y^k\big)^2\bigr)
 + \big(p_2-c y^2\big)\,\d y^2 + \cdots + \big(p_k-c y^k\big)\,\d y^k
\]
and observing that
\begin{gather*}
H\big(y^1 + \tfrac12c \big(y^2\big)^2 + \cdots + \tfrac12c \big(y^k\big)^2\big)\\
\qquad = H\big(y^1\big) + c \big(\d y^2\big)^2 + \cdots + c \big(\d y^k\big)^2
 + c y^2 H\big(y^2\big) + \cdots + c y^k H\big(y^k\big)
\end{gather*}
shows that, setting
\[
\bar y^1 = y^1 + \tfrac12c \big(y^2\big)^2 + \cdots + \tfrac12c \big(y^k\big)^2,
\qquad
\bar y^i = y^i,
\qquad\text{and}\qquad
\bar p_i = p_i - c y^i,
\]
one has
$\omega = \d\bar y^1
 + \bar p_2\,\d \bar y^2 + \cdots + \bar p_k\,\d \bar y^k$.

Thus, one could have chosen the functions~$y^1,\ldots,y^k$, $p_2,\ldots,p_k$
with~$H\big(y^1\big)$ being positive definite on the hyperplane~$K_m$.
Assume now that this was done.

It still needs to be shown that one can choose
the functions~$y^1,\ldots,y^k$, $p_2,\ldots,p_k$
with $H\big(y^1\big)$ being positive definite on all of $T_mM$,
not just on $K_m$, which is the kernel of $\d y^1$ at~$m$.
To do this, note that, if~$\phi$ is any smooth function on a neighborhood
of the origin in~$\bbR$, then
\[
H\bigl(\phi\big(y^1\big)\bigr) = \phi'\big(y^1\big) H\big(y^1\big) + \phi''\big(y^1\big) \big(\d y^1\big)^2.
\]
Hence, by choosing a~$\phi$ with $\phi(0)=0$, $\phi'(0)=1$
and $\phi''(0)>0$ sufficiently large,
I can arrange that~$\phi\big(y^1\big)$ be strictly $\nabla$-convex at~$m$.
Since
\[
\omega = \frac{1}{\phi'(y_1)} \big(\d\bigl(\phi\big(y^1\big)\bigr)
    + \phi'(y_1)p_2\,\d y^2 +\cdots + \phi'(y_1)p_k\,\d y^k\big),
\]
one sees that the functions
$\bar y^1,\ldots,\bar y^k$, $\bar p_2,\ldots,\bar p_k$, where
\[
\bar y^1 = \phi\big(y^1\big),
\qquad\text{and}\qquad
\bar y^i = y^i,
\qquad
\bar p_i = \phi'(y_1)p_i,\qquad 2\le i\le k,
\]
(with $a = 1/\phi'(y_1) >0$), give a Pfaff--Darboux representation
for~$\omega$ that is compatible with the foliation~$\cL$ and
for which $\bar y^1$ is strictly $\nabla$-convex.\footnote{This is the same idea that Ekeland and Nirenberg \cite{MR2023129}
used in their generalization of their Theorem~1
to cover the quasi-convex case.}

Thus, one can assume henceforth that,
on an open $m$-neighborhood~$V_1\subset M$,
one has a~representation of the form
\[
\omega = \d y^1 + p_2\,\d y^2 +\cdots + p_k\,\d y^k,
\]
where the functions~$y^1,\ldots,y^k,p_2,\ldots,p_k\in C^\infty(V_1)$
all vanish at~$m$,
the equations $\d y^i=0$ define the tangents to the leaves of~$\cL$ in~$V_1$,
and~$y^1$ is strictly~$\nabla$-convex.

Under these assumptions, there is a constant~$b>0$ sufficiently large
so that $H\big(y^i+by^1\big) = H\big(y^i\big)+b H\big(y^1\big)$ is positive definite at~$m$
for~$2\le i\le k$.
Thus, writing
\[
\omega = \bigl(1-b(p_2{+}\cdots{+}p_k)\bigr)\,\d y^1 + p_2\,\d\big(y^2+b y^1\big)
   + \cdots + p_k\,\d\big(y^k+b y^1\big),
\]
it follows that I can,
after restricting to an $m$-neighborhood $V_2\subset V_1$ on which
the function $a = \bigl(1-b(p_2{+}\cdots{+}p_k)\bigr)$ is positive,
dividing by~$a>0$,
and replacing~$y^j$ by $y^j+by^1$ and~$p_j$ by~$p_j/a$ for $2\le j\le k$,
assume that I have a representation
\[
\omega = \d y^1 + p_2\,\d y^2 +\cdots + p_k\,\d y^k,
\]
in which all of the~$H\big(y^j\big)$ are positive definite at~$m$,
i.e., the $y^j$ are strictly $\nabla$-convex on some neighborhood of~$m$
and the $p_i$ all vanish at~$m$.

Finally, for $\varepsilon>0$ and sufficiently small, write
\[
\omega = \d\bigl(y^1-\varepsilon\big(y^2{+}\cdots{+}y^k\big)\bigr)
   + (p_2 + \varepsilon)\,\d y^2
   + \cdots + (p_k + \varepsilon)\,\d y^k.
\]
Then, setting~$u^1 = y^1-\varepsilon\big(y^2+\cdots+y^k\big)$ and~$u^j = y^j$
for~$j>1$ and setting~$a_1 = 1$ and~${a_j = \varepsilon + p_j}$ for~$j>1$,
one achieves the desired convex Pfaff--Darboux representation
on an open $m$-neigh\-bor\-hood~$U\subset V_2$.
\end{proof}
